\section{一次完整的前向与反向传播}
\label{sec:14-Deep-Learning/02-Losses-and-Backpropagation/02}

一个自己手写的打分层上线到训练脚本里，损失曲线看着正常，往下掉得挺稳。麻烦出在换批大小之后：从 \(32\) 改到 \(256\)，学习率不动就发散，改小一个数量级又收敛得极慢。同事的判断是``梯度肯定没写错，损失在降''。这句话不成立——一个方向错了的梯度也可能让损失下降，只是降的不是你以为的那个目标。要把这类问题定位下来，得先知道一次反向传播里每个数是怎么来的。

上一节确定了损失，也就确定了``什么算好''。剩下的问题是把这个标量的偏好分配到十万个参数上：每个参数各自该往哪个方向、迈多大一步。逐参数用有限差分近似的话，一条导数一次完整前向，十万个导数就是十万次前向，一个 epoch 都跑不完。

反向传播换了个走法：利用计算图的局部结构，一趟前向加一趟反向，把所有参数的梯度一次拿全，总代价与一次前向同阶。它只负责算梯度，不负责怎么用梯度——梯度下降、Adam 或别的优化器是后面的事。这条``代价与前向同阶''的结论以及一般计算图上的推导，在\hyperref[sec:10-Matrix-Calculus/04-Applications/03]{``计算图上的反向传播与自动微分''}里有完整交代；这一节只做一件事，把一次数值过程从头走到尾，再说清楚出问题时该查哪里。

\subsection{先在一个标量网络上把每个数走一遍}

考虑一个最小网络：

\[
z=w_1x+b_1,\qquad
h=\relu(z),\qquad
o=w_2h+b_2,
\]

\[
L=\frac12(o-y)^2.
\]

取一组具体数值：\(x=2\)，\(w_1=0.5\)，\(b_1=0\)，\(w_2=-1\)，\(b_2=0\)，标签 \(y=1\)。前向从输入流向损失：

\[
z=0.5\times 2+0=1,\quad
h=\max(0,1)=1,\quad
o=(-1)\times 1+0=-1,\quad
L=\frac12(-1-1)^2=2.
\]

前向途中每算出一个中间值就存下来：\(z\)、\(h\)、\(o\)。反向要用到它们，这也是训练比推理吃显存得多的原因——存的不是参数，是激活值。

\begin{center}
\begin{tikzpicture}[>=stealth,
  unit/.style={draw,rounded corners=2pt,minimum width=1.25cm,minimum height=0.7cm,fill=blue!6},
  every node/.style={font=\small}]
  \node[unit] (x) at (-5.2,0) {\(x\)};
  \node[unit] (z) at (-2.6,0) {\(z=w_1x+b_1\)};
  \node[unit] (h) at (0,0) {\(h=\relu(z)\)};
  \node[unit] (o) at (2.6,0) {\(o=w_2h+b_2\)};
  \node[unit] (l) at (5.2,0) {\(L\)};
  \draw[->,blue!70!black,thick] (x) -- (z);
  \draw[->,blue!70!black,thick] (z) -- (h);
  \draw[->,blue!70!black,thick] (h) -- (o);
  \draw[->,blue!70!black,thick] (o) -- (l);
  \draw[<-,black,thick] (x.south) -- ++(0,-0.65);
  \draw[->,black,thick] (l.south) -- ++(0,-0.65);
  \draw[<-,black,thick] (z.south) -- ++(0,-0.65);
  \draw[<-,black,thick] (h.south) -- ++(0,-0.65);
  \draw[<-,black,thick] (o.south) -- ++(0,-0.65);
  \draw[->,black,thick] (5.2,-1.0) -- (-5.2,-1.0);
  \node[above] at (0,0.55) {前向：计算并保存中间值};
  \node[below] at (0,-1.0) {反向：上游梯度 \(\times\) 局部导数};
\end{tikzpicture}

\emph{前向箭头携带数值，反向箭头携带敏感度；每经过一个节点，就把上游敏感度乘以该节点的局部导数。}
\end{center}

现在照着图里下面那条箭头往回走。\(L=\frac12(o-y)^2\) 对 \(o\) 的导数是直接的：

\[
\frac{\partial L}{\partial o}=o-y=-1-1=-2 .
\]

这个 \(-2\) 读作：输出 \(o\) 每增加一个微小量，损失大约变化 \(-2\) 倍于那个量。想让损失下降，\(o\) 该往正方向移动。

接着过 \(o=w_2h+b_2\) 这道门。\(h\) 在前向时存下来是 \(1\)，于是

\[
\frac{\partial L}{\partial w_2}
=\frac{\partial L}{\partial o}\cdot\frac{\partial o}{\partial w_2}
=(-2)\times h=-2,
\qquad
\frac{\partial L}{\partial b_2}=(-2)\times 1=-2 .
\]

再穿过 \(h=\relu(z)\)。ReLU 的导数在 \(z>0\) 时为 \(1\)、\(z<0\) 时为 \(0\)，前向记下 \(z=1>0\)，所以这道门是通的。先看敏感度怎么传到 \(h\)：

\[
\frac{\partial L}{\partial h}
=\frac{\partial L}{\partial o}\cdot\frac{\partial o}{\partial h}
=(-2)\times w_2=2 .
\]

这个 \(2\) 是正的，意思是 \(h\) 变大损失会变大——因为 \(w_2\) 为负，\(h\) 增大把 \(o\) 推得更负，离 \(y=1\) 更远。继续退到 \(z\)：

\[
\frac{\partial L}{\partial z}
=\frac{\partial L}{\partial h}\cdot \relu'(z)
=2\times 1=2 ,
\]

最后到第一层参数：

\[
\frac{\partial L}{\partial w_1}
=\frac{\partial L}{\partial z}\cdot x
=2\times 2=4,\qquad
\frac{\partial L}{\partial b_1}=2 .
\]

四个梯度到手：\(w_2\) 是 \(-2\)，\(b_2\) 是 \(-2\)，\(w_1\) 是 \(4\)，\(b_1\) 是 \(2\)。一趟前向、一趟反向，没有任何一段计算被重复。

整个过程里反复出现的动作只有两个。一是把上游传来的敏感度乘上当前操作的局部导数，这就是链式法则按从右往左的顺序乘；二是当一个变量通过多条路径影响损失时，各条路径的贡献要相加。第二条在这个例子里没派上用场，\(z\)、\(h\)、\(o\) 各自只有一条通往损失的线路。网络一旦出现分支——跳跃连接、共享权重、多任务头——汇合点上的加法就是必需的，漏掉它梯度会小得毫无规律。

\subsection{批量矩阵形式只是同一套动作的压缩写法}

真实训练不会一条样本一个标量地算。把批内样本压进矩阵之后，公式变的是排布，不是逻辑。

采用样本按行的约定。批大小 \(B\)，各层宽度 \(d_0,d_1,d_2\)：

\[
X\in\R^{B\times d_0},\quad
W_1\in\R^{d_0\times d_1},\quad
W_2\in\R^{d_1\times d_2},
\]
\[
Z_1\in\R^{B\times d_1},\quad
H\in\R^{B\times d_1},\quad
Z_2\in\R^{B\times d_2}.
\]

前向是

\[
Z_1=XW_1+\mathbf1b_1^{\trans},
\qquad H=\phi(Z_1),
\]
\[
Z_2=HW_2+\mathbf1b_2^{\trans},
\qquad
L=\frac1B\sum_{i=1}^B\ell(Z_{2i},y_i).
\]

反向从损失对各样本 logits 的导数矩阵起步，记

\[
G_2=\frac{\partial L}{\partial Z_2}\in\R^{B\times d_2}.
\]

逐样本看就清楚了：\(Z_2\) 的第 \(i\) 行只进入第 \(i\) 个样本的损失，它对 \(W_2\) 梯度的贡献是隐藏行向量 \(h_i\)（长度 \(d_1\)）与敏感度行向量 \(g_{2i}\)（长度 \(d_2\)）的外积，是一个 \(d_1\times d_2\) 矩阵。对 \(i\) 求和恰好就是矩阵乘法 \(H^{\trans}G_2\)。于是

\[
\nabla_{W_2}L=H^{\trans} G_2,\qquad
\nabla_{b_2}L=G_2^{\trans}\mathbf1 .
\]

偏置在前向里被广播到每一行，反向就要沿批轴求和，把各样本的贡献收回到同一个向量上。梯度继续跨过 \(Z_2=HW_2\) 传到 \(H\)，得到 \(G_2W_2^{\trans}\)，样本维保持不变而隐藏维从 \(d_2\) 变回 \(d_1\)；再逐元素乘上激活函数的导数：

\[
G_1
=(G_2W_2^{\trans})\odot\phi'(Z_1).
\]

第一层参数梯度同理，

\[
\nabla_{W_1}L=X^{\trans} G_1,\qquad
\nabla_{b_1}L=G_1^{\trans}\mathbf1 .
\]

对 \(\softmax\) 交叉熵这个最常见的组合，\(G_2\) 有一个干净的闭式：

\[
G_2=\frac{P-Y}{B},
\]

其中 \(P\) 是 \(Z_2\) 逐行 \(\softmax\) 之后的概率矩阵，\(Y\) 是 one-hot 标签矩阵。这就是上一节末尾那个 \(p-y\) 信号，逐样本逐类别算完再除以批大小。代进上面几条公式，两层分类器的全部参数梯度就齐了。

框架里的反向模式自动微分在这个两层网络和一个千层模型上执行的是同一种运算：把上游梯度与局部 Jacobian 做向量--Jacobian 乘积，沿计算图反向推进。区别只在图更大、节点更多，以及图结构和缓存由框架替你维护。

\subsection{梯度出问题时先查形状和因子，再查数值}

回到开头那次换批大小的事故。矩阵公式里藏着一个极易被掩盖的因子：如果损失按批取平均，那个 \(1/B\) 必须已经出现在 \(G_2\) 里。漏掉它，梯度整体大 \(B\) 倍，效果等同于学习率被悄悄放大 \(B\) 倍——批大小从 \(32\) 改到 \(256\)，等效学习率跟着涨八倍，发散得理所当然。反过来在自定义损失里又除了一次，梯度被压到几乎不动，看起来像``模型学不动''。换批大小之后学习率忽然不对，第一站就该查这个因子；掩码和样本权重的分母也属于同一类问题。

比数值更早能查的是形状。每个参数梯度的形状必须与参数本身逐维相同，这条断言写进代码只要一行，却能挡住转置方向写反、批轴规约方向搞错这两类最常见的手写错误。形状对上之后再看量级：各层梯度范数如果相差好几个数量级，问题多半在初始化或归一化，而不在反向公式。

真正难查的是那些不改变形状、也不报错的错误。原地修改了反向还要用的缓存张量，无意中把某段计算从图上摘了下来，网络与损失两边的转置约定不一致，多路径贡献互相覆盖而不是累加——这几种都会让梯度变成一个与正确值无关的张量，而损失照样下降。这时候需要一个独立的验证手段。

\subsection{梯度检查是局部证据，不是全局担保}

有限差分提供了这样一个手段。沿一个随机方向 \(v\) 做中心差分，

\[
\frac{L(\theta+\varepsilon v)-L(\theta-\varepsilon v)}{2\varepsilon}
\approx
\nabla L(\theta)^{\trans} v .
\]

中心差分里 \(O(\varepsilon^2)\) 的 Hessian 项相消，截断误差比单侧差分小一个量级。不必逐参数核对——转置或规约类的错误几乎一定会破坏梯度与随机方向的内积关系，取两三个随机方向各查一次，漏检的概率已经很低。步长 \(\varepsilon\) 要在截断误差和舍入误差之间折中，双精度下取 \(10^{-5}\) 到 \(10^{-6}\) 通常合适，单精度下这个检查本身就不太可靠。

通过了梯度检查，说明的仅仅是：在当前这个参数点附近，偏导数的数值实现大致正确。它不说明模型选对了，不说明损失与标签的语义匹配，不说明数据切分干净，更不说明训练会收敛。它还有一处已知的盲区：ReLU 在零点左右导数不等，\(z\) 落在零附近的样本会让有限差分给出不可信的结果，检查时应当避开这些点，或者换一个光滑激活来做验证。

梯度算对之后，训练仍然可能失败——步长把参数推出有效区域，深层网络的梯度在连乘中消失或爆炸，批噪声让轨迹在最优点附近来回打转。这些问题不在计算图上，而在优化过程本身，下一单元接手。
